\chapter{Building a CLI Tool}\label{ch:building-cli-tool}

\index{CLI tool}
\index{command-line application}

Every language earns its stripes when you build something real with it.
So far we have explored Lattice's features one by one---phases, concurrency, pattern matching, modules.
Now it is time to weave them together into a complete, polished command-line application: a file-search utility called \texttt{find-it} that searches files for a pattern and prints matching lines with context.

Along the way we will use three libraries from the Lattice ecosystem:
\lstinline|cli| for argument parsing,
\lstinline|dotenv| for loading configuration from \texttt{.env} files,
and \lstinline|log| for structured, leveled logging.
By the end of this chapter you will have a template for every CLI tool you write in Lattice.

%% ─────────────────────────────────────────────────────────────────────────────
\section{The \texttt{cli} Library for Argument Parsing}\label{sec:cli-library}
%% ─────────────────────────────────────────────────────────────────────────────

\index{cli library}
\index{argument parsing}

A command-line tool lives or dies by its interface. Users expect \texttt{-{}-help} to work, short flags like \texttt{-v} to combine, and helpful error messages when they get something wrong.
The \lstinline|cli| library (found in \filename{lib/cli.lat}) gives us all of this through a declarative API.

\subsection{Creating a Parser}\label{sec:cli-creating-parser}

We begin by importing the library and creating a new application definition:

\begin{lstlisting}[language=Lattice, caption={A minimal CLI parser}]
import "lib/cli" as cli

let app = cli.new("greet", "A friendly greeter")
app.flag("loud", "l", "Shout the greeting")
app.option("name", "n", "Who to greet", "World")

let opts = app.parse()
// opts is a Map with keys: "loud", "name", "_args"
\end{lstlisting}

The \lstinline|cli.new| function takes two arguments---the program name and a short description---and returns a Map packed with closures for configuring and executing the parser.

\begin{defnbox}{CLI Application Map}
The value returned by \lstinline|cli.new| is a Map containing closures:
\lstinline|.flag(name, short, desc)| registers a boolean flag,
\lstinline|.option(name, short, desc, default)| registers a named option with a default value,
\lstinline|.arg(name, desc, required)| registers a positional argument,
\lstinline|.help()| generates help text, and
\lstinline|.parse()| parses the command-line arguments and returns a result Map.
\end{defnbox}

\subsection{Flags, Options, and Positional Arguments}\label{sec:cli-flag-option-arg}

\index{flags!CLI}
\index{options!CLI}
\index{positional arguments}

The \lstinline|cli| library distinguishes three kinds of command-line input:

\begin{itemize}
  \item \textbf{Flags} are boolean switches. They default to \lstinline|false| and become \lstinline|true| when present. Think \texttt{-{}-verbose} or \texttt{-v}.
  \item \textbf{Options} carry a value. They can appear as \texttt{-{}-output file.txt}, \texttt{-{}-output=file.txt}, or \texttt{-o file.txt}. Every option has a default.
  \item \textbf{Positional arguments} are the bare values that remain after flags and options are consumed. They can be required or optional.
\end{itemize}

\begin{lstlisting}[language=Lattice, caption={Registering all three kinds of input}]
import "lib/cli" as cli

let app = cli.new("findi", "Search files for a pattern")

// Flags (boolean)
app.flag("verbose", "v", "Show extra information")
app.flag("count", "c", "Print only the count of matches")
app.flag("ignorecase", "i", "Case-insensitive search")

// Options (with default values)
app.option("output", "o", "Write results to a file", "")
app.option("context", "C", "Lines of context to show", "0")

// Positional arguments
app.arg("pattern", "The search pattern", true)     // required
app.arg("path", "File or directory to search", false) // optional
\end{lstlisting}

\subsection{Parsing and the Result Map}\label{sec:cli-parsing}

Calling \lstinline|app.parse()| reads the process arguments (via the built-in \lstinline|args()| function), matches them against the registered flags, options, and positionals, and returns a result Map.
Flags map to \lstinline|Bool| values, options map to \lstinline|String| values (or their defaults), and positional arguments map to \lstinline|String| values.
Any extra positional arguments land in the \lstinline|"_args"| key as an array.

\begin{lstlisting}[language=Lattice, caption={Working with parsed results}]
let opts = app.parse()

// Boolean flags
let is_verbose = opts.get("verbose")    // true or false

// String options (always strings---convert if needed)
let ctx_lines = to_int(opts.get("context"))

// Positional arguments
let pattern = opts.get("pattern")       // guaranteed present (required)
let path = opts.get("path")             // may be nil if not provided
\end{lstlisting}

\begin{warnbox}{Options Are Always Strings}
The \lstinline|cli| library returns all option values as strings, including numeric defaults.
When you register \lstinline|app.option("port", "p", "Port number", "8080")|, the parsed value will be the string \lstinline|"8080"|, not the integer \lstinline|8080|.
Use \lstinline|to_int()| or \lstinline|to_float()| to convert.
\end{warnbox}

\subsection{Automatic Help and Error Handling}\label{sec:cli-help-errors}

\index{help flag}

The parser automatically intercepts \texttt{-{}-help} and \texttt{-h}, printing a nicely formatted help message and exiting.
You never need to register a help flag yourself.

Running our \texttt{findi} tool with \texttt{-{}-help} produces:

\begin{lstlisting}[language={}, caption={Auto-generated help output}]
findi - Search files for a pattern

Usage: findi [OPTIONS] <pattern> [path]

Arguments:
  <pattern>     The search pattern (required)
  [path]        File or directory to search

Options:
  -v, --verbose       Show extra information
  -c, --count         Print only the count of matches
  -i, --ignorecase    Case-insensitive search
  -o, --output        Write results to a file [default: ]
  -C, --context       Lines of context to show [default: 0]
  -h, --help          Show this help message
\end{lstlisting}

If a user provides an unknown flag or forgets a required argument, the parser prints an error along with the help text and exits with code 1.

\subsection{Combined Short Flags}\label{sec:cli-combined-flags}

\index{combined flags}

Like many Unix tools, the \lstinline|cli| library supports combining short flags into a single token.
The invocation \texttt{findi -vic "TODO" src/} is equivalent to \texttt{findi -v -i -c "TODO" src/}.
When combining, every character except the last must be a flag.
The last character may be an option, in which case the next argument becomes its value:

\begin{lstlisting}[language={}, caption={Combined flags with a trailing option}]
# -v and -i are flags; -C is an option consuming "3"
findi -viC 3 "pattern" src/
\end{lstlisting}

\begin{tipbox}{Double-Dash Separator}
The \lstinline|cli| library supports the \texttt{-{}-} separator.
Everything after \texttt{-{}-} is treated as a positional argument, even if it starts with a dash.
This is essential for patterns like \texttt{findi -{}- -TODO- file.txt} where the pattern itself looks like a flag.
\end{tipbox}

%% ─────────────────────────────────────────────────────────────────────────────
\section{The \texttt{dotenv} Library for Configuration}\label{sec:dotenv-library}
%% ─────────────────────────────────────────────────────────────────────────────

\index{dotenv library}
\index{environment variables}
\index{.env file}

CLI tools often need configuration that varies between environments---database URLs, API keys, default directories.
Hardcoding these into source files is fragile.
The \lstinline|dotenv| library (in \filename{lib/dotenv.lat}) loads key-value pairs from \texttt{.env} files into the process environment, so your code can read them with \lstinline|env()|.

\subsection{The Basics: \texttt{load} and \texttt{env}}\label{sec:dotenv-basics}

Suppose our search tool needs a default search directory.
We create a \texttt{.env} file alongside our script:

\begin{lstlisting}[language={}, caption={A \texttt{.env} file}]
# Default configuration for findi
FINDI_DEFAULT_PATH=./src
FINDI_MAX_RESULTS=100
FINDI_LOG_LEVEL=info
\end{lstlisting}

Then we load it at startup:

\begin{lstlisting}[language=Lattice, caption={Loading a .env file}]
import "lib/dotenv" as dotenv

dotenv.load()  // loads .env from the current directory

let default_path = env("FINDI_DEFAULT_PATH")  // "./src"
let max_results = to_int(env("FINDI_MAX_RESULTS"))  // 100
\end{lstlisting}

The \lstinline|dotenv.load()| function is deliberately forgiving: if no \texttt{.env} file exists, it silently does nothing.
This is ideal for production, where configuration typically comes from the real environment rather than a file.

\subsection{Loading Specific Files}\label{sec:dotenv-load-file}

\index{dotenv!load\_file}

When you need to load from a particular path---say, a per-environment configuration---use \lstinline|dotenv.load_file|:

\begin{lstlisting}[language=Lattice, caption={Loading a specific .env file}]
import "lib/dotenv" as dotenv

dotenv.load_file("config/.env.production")
\end{lstlisting}

Unlike \lstinline|load()|, this function will halt with an error if the file does not exist.
It is meant for situations where the configuration file is mandatory.

\subsection{Advanced Loading with Options}\label{sec:dotenv-load-opts}

\index{dotenv!load\_opts}

For more control, \lstinline|dotenv.load_opts| accepts a Map of options:

\begin{lstlisting}[language=Lattice, caption={Loading with options}]
import "lib/dotenv" as dotenv

let opts = Map::new()
opts.set("path", ".env.local")
opts.set("override", true)           // overwrite existing env vars
opts.set("required", ["DATABASE_URL", "SECRET_KEY"])

dotenv.load_opts(opts)
\end{lstlisting}

\begin{notebox}{Override Behavior}
By default, \lstinline|dotenv| never overwrites environment variables that are already set.
This means a real \lstinline|DATABASE_URL| set in your shell takes precedence over whatever is in \texttt{.env}.
Set \lstinline|"override"| to \lstinline|true| only when you explicitly want the file to win---typically during local development.
\end{notebox}

The \lstinline|"required"| key is particularly useful for fail-fast behavior.
If any listed variable is missing after loading, the program halts with a clear error message rather than silently using \lstinline|nil| values.

\subsection{Parsing Rules}\label{sec:dotenv-parsing-rules}

\index{dotenv!parsing rules}

The parser in \filename{lib/dotenv.lat} handles a rich set of \texttt{.env} formats.
Here is a file that exercises most of them:

\begin{lstlisting}[language={}, caption={A feature-rich .env file}]
# Comments start with #
BASIC=value

# Whitespace around = is trimmed
SPACED = hello

# Double-quoted values support escapes
GREETING="Hello,\nWorld!"

# Single-quoted values are literal (no escapes)
RAW='Hello,\nWorld!'

# Inline comments on unquoted values
PORT=3000 # default port

# export prefix is supported
export API_KEY=sk-abc123

# Variable expansion in double-quoted values
BASE_URL="https://api.example.com"
FULL_URL="${BASE_URL}/v2"

# Multiline double-quoted values
CERTIFICATE="-----BEGIN CERT-----
MIIB...
-----END CERT-----"
\end{lstlisting}

\begin{tipbox}{Parsing Without Setting}
Sometimes you want to inspect a \texttt{.env} file without modifying the environment.
Use \lstinline|dotenv.parse(".env")| to get a Map of key-value pairs, or \lstinline|dotenv.parse_string(content)| to parse an in-memory string.
\end{tipbox}

%% ─────────────────────────────────────────────────────────────────────────────
\section{The \texttt{log} Library for Structured Logging}\label{sec:log-library}
%% ─────────────────────────────────────────────────────────────────────────────

\index{log library}
\index{logging}
\index{structured logging}

A \lstinline|print| statement is fine for quick debugging, but production CLI tools need leveled, timestamped, and optionally structured logging.
The \lstinline|log| library (in \filename{lib/log.lat}) provides exactly that.

\subsection{Quick Start: Module-Level Functions}\label{sec:log-quick-start}

The fastest way to start logging is with the module-level convenience functions:

\begin{lstlisting}[language=Lattice, caption={Module-level logging}]
import "lib/log" as log

log.info("Server started on port 8080")
log.warn("Connection pool running low")
log.error("Failed to read configuration file")
log.debug("Cache hit ratio: 0.87")  // suppressed at default level
\end{lstlisting}

These write to \texttt{stderr} using the default logger, which has a minimum level of \lstinline|info|.
That means \lstinline|log.debug| messages are silenced unless you create a custom logger.

The output is human-readable, timestamped text:

\begin{lstlisting}[language={}, caption={Default log output format}]
2025-01-15 10:30:45 [INFO] Server started on port 8080
2025-01-15 10:30:45 [WARN] Connection pool running low
2025-01-15 10:30:45 [ERROR] Failed to read configuration file
\end{lstlisting}

\subsection{Log Levels}\label{sec:log-levels}

\index{log levels}

The library defines four levels in increasing order of severity:

\begin{center}
\begin{tabular}{llp{7cm}}
\toprule
\textbf{Level} & \textbf{Value} & \textbf{Purpose} \\
\midrule
\lstinline|debug| & 0 & Detailed diagnostic information for development \\
\lstinline|info|  & 1 & Normal operational messages \\
\lstinline|warn|  & 2 & Something unexpected happened, but the tool can continue \\
\lstinline|error| & 3 & Something failed; the operation cannot proceed normally \\
\bottomrule
\end{tabular}
\end{center}

A logger only emits messages at or above its configured minimum level.
Set the level to \lstinline|"debug"| during development and \lstinline|"warn"| or \lstinline|"error"| in quiet production modes.

\subsection{Creating a Custom Logger}\label{sec:log-custom-logger}

\index{log!new}

For full control, create a logger with \lstinline|log.new|:

\begin{lstlisting}[language=Lattice, caption={Creating a custom logger}]
import "lib/log" as log

let cfg = Map::new()
cfg.set("level", "debug")           // show everything
cfg.set("file", "findi.log")        // also write to a file
cfg.set("format", "json")           // structured JSON output
cfg.set("prefix", "[findi]")        // prefix every message

let logger = log.new(cfg)

logger.info("Starting search")
logger.debug("Scanning directory: ./src")
\end{lstlisting}

The configuration Map supports four keys:

\begin{itemize}
  \item \lstinline|"level"| --- minimum level: \lstinline|"debug"|, \lstinline|"info"|, \lstinline|"warn"|, or \lstinline|"error"| (default: \lstinline|"info"|).
  \item \lstinline|"file"| --- path to a log file. Messages are appended to this file \emph{in addition to} stderr (default: none).
  \item \lstinline|"format"| --- \lstinline|"text"| for human-readable lines, \lstinline|"json"| for machine-readable JSON (default: \lstinline|"text"|).
  \item \lstinline|"prefix"| --- a string prepended to every message (default: \lstinline|""|).
\end{itemize}

\subsection{Structured Context}\label{sec:log-context}

\index{log!context}

Each logging method accepts an optional second argument: a Map of contextual data.
In text mode, the context is serialized as JSON and appended to the line.
In JSON mode, the context keys are merged into the top-level log entry:

\begin{lstlisting}[language=Lattice, caption={Logging with context}]
let ctx = Map::new()
ctx.set("path", "/api/users")
ctx.set("status", 404)
ctx.set("elapsed_ms", 12)

logger.warn("Route not found", ctx)
\end{lstlisting}

In text format, this produces:

\begin{lstlisting}[language={}, caption={Text log with context}]
2025-01-15 10:30:47 [WARN] [findi] Route not found {"path":"/api/users","status":404,"elapsed_ms":12}
\end{lstlisting}

In JSON format:

\begin{lstlisting}[language={}, caption={JSON log with context}]
{"timestamp":"2025-01-15T10:30:47","level":"WARN","prefix":"[findi]","message":"Route not found","path":"/api/users","status":404,"elapsed_ms":12}
\end{lstlisting}

\begin{tipbox}{Changing Levels at Runtime}
Because Lattice closures capture their environment by value, you cannot mutate a logger's level in place.
Instead, use \lstinline|log.set_level(logger, "error")| to create a \emph{new} logger with the same configuration but a different threshold.
This is an intentional design choice---loggers are value-like and safe to share.
\end{tipbox}

%% ─────────────────────────────────────────────────────────────────────────────
\section{Putting It All Together: A Complete CLI Application}\label{sec:cli-complete-app}
%% ─────────────────────────────────────────────────────────────────────────────

\index{findi@\texttt{findi} (example application)}

Now let us build \texttt{findi}, our file-search tool.
It reads configuration from \texttt{.env}, parses arguments with \lstinline|cli|, and logs its activity with \lstinline|log|.
The tool searches a directory tree for lines matching a pattern and prints the results.

\subsection{Project Layout}\label{sec:cli-project-layout}

\begin{lstlisting}[language={}, caption={Project structure}]
findi/
  findi.lat          # main entry point
  .env               # default configuration
\end{lstlisting}

And our \texttt{.env} file:

\begin{lstlisting}[language={}, caption={Default .env configuration}]
FINDI_DEFAULT_PATH=.
FINDI_LOG_LEVEL=info
FINDI_LOG_FORMAT=text
\end{lstlisting}

\subsection{The Complete Source}\label{sec:cli-complete-source}

\begin{lstlisting}[language=Lattice, caption={findi.lat --- A complete CLI file-search tool}]
import "lib/cli" as cli
import "lib/dotenv" as dotenv
import "lib/log" as log

// ── Load environment configuration ──────────────────────
dotenv.load()

// ── Set up argument parser ──────────────────────────────
let app = cli.new("findi", "Search files for a text pattern")

app.flag("verbose", "v", "Show verbose output")
app.flag("count", "c", "Print only the match count")
app.flag("ignorecase", "i", "Case-insensitive matching")
app.option("output", "o", "Write results to a file", "")
app.option("context", "C", "Lines of context around matches", "0")
app.arg("pattern", "The text pattern to search for", true)
app.arg("path", "File or directory to search", false)

let opts = app.parse()

// ── Configure logging ───────────────────────────────────
let log_cfg = Map::new()
let env_level = env("FINDI_LOG_LEVEL")
if typeof(env_level) == "String" {
    log_cfg.set("level", env_level)
} else {
    log_cfg.set("level", "info")
}
let env_format = env("FINDI_LOG_FORMAT")
if typeof(env_format) == "String" {
    log_cfg.set("format", env_format)
}
if opts.get("verbose") {
    log_cfg.set("level", "debug")
}
let logger = log.new(log_cfg)

// ── Extract parsed arguments ────────────────────────────
let pattern = opts.get("pattern")
let ignore_case = opts.get("ignorecase")
let count_only = opts.get("count")
let ctx_lines = to_int(opts.get("context"))
let output_file = opts.get("output")

// Determine search path: CLI arg > env default > "."
flux search_path = "."
let arg_path = opts.get("path")
if arg_path != nil {
    search_path = arg_path
} else {
    let env_path = env("FINDI_DEFAULT_PATH")
    if typeof(env_path) == "String" {
        search_path = env_path
    }
}

logger.debug("Configuration loaded")
let debug_ctx = Map::new()
debug_ctx.set("pattern", pattern)
debug_ctx.set("path", search_path)
debug_ctx.set("ignore_case", ignore_case)
logger.debug("Search parameters", debug_ctx)

// ── Search logic ────────────────────────────────────────

fn search_file(file_path: String, search_pattern: String,
               case_insensitive: Bool) -> Array {
    let matches = []
    let content = try {
        read_file(file_path)
    } catch e {
        return []
    }
    let lines = content.split("\n")
    flux line_num = 1
    for line in lines {
        flux haystack = line
        flux needle = search_pattern
        if case_insensitive {
            haystack = line.to_lower()
            needle = search_pattern.to_lower()
        }
        if haystack.contains(needle) {
            let hit = Map::new()
            hit.set("file", file_path)
            hit.set("line_num", line_num)
            hit.set("text", line)
            matches.push(hit)
        }
        line_num = line_num + 1
    }
    return matches
}

fn collect_files(dir_path: String) -> Array {
    let entries = list_dir(dir_path)
    flux files = []
    for entry in entries {
        let full_path = dir_path + "/" + entry
        if is_dir(full_path) {
            let sub_files = collect_files(full_path)
            for sf in sub_files {
                files.push(sf)
            }
        } else {
            files.push(full_path)
        }
    }
    return files
}

// ── Execute the search ──────────────────────────────────
flux all_matches = []

if is_dir(search_path) {
    let files = collect_files(search_path)
    logger.info("Scanning " + to_string(len(files)) + " files")
    for file_path in files {
        let hits = search_file(file_path, pattern, ignore_case)
        for hit in hits {
            all_matches.push(hit)
        }
    }
} else {
    logger.info("Searching single file: " + search_path)
    all_matches = search_file(search_path, pattern, ignore_case)
}

// ── Format and output results ───────────────────────────
if count_only {
    print(to_string(len(all_matches)))
} else {
    flux output_text = ""
    for hit in all_matches {
        let line = hit.get("file") + ":"
            + to_string(hit.get("line_num")) + ":"
            + hit.get("text")
        output_text = output_text + line + "\n"
    }
    if output_file != "" {
        write_file(output_file, output_text)
        logger.info("Results written to " + output_file)
    } else {
        print(output_text)
    }
}

let result_ctx = Map::new()
result_ctx.set("matches", len(all_matches))
result_ctx.set("path", search_path)
logger.info("Search complete", result_ctx)
\end{lstlisting}

\subsection{Running the Tool}\label{sec:cli-running}

Let us take our tool for a spin:

\begin{lstlisting}[language={}, caption={Running findi}]
$ lattice findi.lat "import" ./src
./src/main.lat:1:import "lib/cli" as cli
./src/main.lat:2:import "lib/dotenv" as dotenv
./src/utils.lat:1:import "lib/log" as log

$ lattice findi.lat -c "TODO" ./src
7

$ lattice findi.lat -vi --context=2 "error" ./src
# verbose logging enabled, case-insensitive,
# with 2 lines of context around each match

$ lattice findi.lat --help
findi - Search files for a text pattern
...
\end{lstlisting}

\subsection{Design Walkthrough}\label{sec:cli-design-walkthrough}

Let us highlight a few design choices worth remembering:

\begin{enumerate}
  \item \textbf{Configuration cascade.} The search path has three sources: the CLI argument (highest priority), the \lstinline|FINDI_DEFAULT_PATH| environment variable, and the hardcoded default \lstinline|"."|.
  This cascade---command line overrides environment overrides defaults---is a pattern you will use in nearly every CLI tool.

  \item \textbf{Verbose flag controls log level.} Rather than building a custom verbosity system, we simply lower the log level to \lstinline|"debug"| when \lstinline|--verbose| is passed.
  All those \lstinline|logger.debug()| calls are free in normal operation.

  \item \textbf{Fail fast with \lstinline|dotenv.load_opts|.} For a tool that requires certain configuration (like a database URL), swap \lstinline|dotenv.load()| for \lstinline|dotenv.load_opts| with a \lstinline|"required"| list.
  The tool will refuse to start if anything is missing, with a clear error message.

  \item \textbf{Separate search logic from I/O.} The \lstinline|search_file| function returns data (an array of match Maps), while the output section at the bottom handles formatting and writing.
  This separation makes it straightforward to add new output formats later.
\end{enumerate}

\begin{notebox}{Why Not Use \texttt{args()} Directly?}
You could parse \lstinline|args()| yourself with string manipulation, and for a two-flag script that might be fine.
But the \lstinline|cli| library handles edge cases you will eventually hit: combined short flags, \texttt{-{}-} separators, missing required arguments, \texttt{-{}-help} generation, and \texttt{-{}-name=value} syntax.
Writing all of that by hand is error-prone and tedious.
Let the library do the work.
\end{notebox}

%% ─────────────────────────────────────────────────────────────────────────────
\section*{Exercises}\label{sec:cli-exercises}
%% ─────────────────────────────────────────────────────────────────────────────

\begin{enumerate}
  \item \textbf{Add a \texttt{-{}-max} option} to \texttt{findi} that limits the number of results printed.
  Default it to \lstinline|0| (unlimited).
  Stop searching early once the limit is reached.

  \item \textbf{Build a \texttt{wordcount} CLI tool} that counts words, lines, and characters in a file (like \texttt{wc}).
  Accept a \texttt{-{}-lines}, \texttt{-{}-words}, and \texttt{-{}-chars} flag to control which counts are shown.
  Use the \lstinline|cli| library for argument parsing and \lstinline|log| for debug output.

  \item \textbf{Environment file chain.} Modify the \texttt{findi} startup to load \texttt{.env} first, then \texttt{.env.local} with \lstinline|override| set to \lstinline|true|.
  This lets developers have personal overrides that are not committed to version control.

  \item \textbf{JSON output mode.} Add a \texttt{-{}-json} flag to \texttt{findi} that outputs results as a JSON array of objects, each with \lstinline|"file"|, \lstinline|"line"|, and \lstinline|"text"| keys.
  Use \lstinline|json_stringify| for the serialization.

  \item \textbf{Custom log format.} Experiment with the \lstinline|log| library's JSON format.
  Create a logger that writes JSON to a file and text to stderr.
  (Hint: you will need two loggers.)
\end{enumerate}

%% ─────────────────────────────────────────────────────────────────────────────
\section*{What's Next}\label{sec:cli-whats-next}
%% ─────────────────────────────────────────────────────────────────────────────

We have built a real command-line tool from scratch, using Lattice's \lstinline|cli|, \lstinline|dotenv|, and \lstinline|log| libraries to handle the plumbing that every serious tool needs.
In the next chapter, we shift from the terminal to the network.
We will build a complete web service---with HTTP routing, HTML templates, a database, and JSON APIs---using Lattice's \lstinline|http_server|, \lstinline|template|, and \lstinline|orm| libraries.
If you thought building a CLI was satisfying, wait until you see your first Lattice web page load in a browser.
